
Reinforcement Learning from Human Feedback (RLHF) has become the dominant paradigm for aligning language models with human preferences \citep{ouyang2022instructgpt}. Current approaches optimize a single objective: helpfulness as judged by human annotators or proxy reward models. This unidirectional alignment ($AI\rightarrow Human)$ has produced capable assistants, yet models follow instructions inconsistently, particularly explicit constraints on format, length, and structure (Zhou et al., 2023a).

Sun et al. (2024) propose a bidirectional alignment framework distinguishing two directions: $AI\rightarrow Human$ (helpfulness, harmlessness) and $Human\rightarrow AI$ (controllability, steerability). While the theoretical framework is compelling, no existing work implements bidirectional training signals. IFEval (Zhou et al., 2023a) measures instruction-following for evaluation; AlpacaEval \citep{dubois2024alpacaeval} measures helpfulness---but these metrics remain siloed for post-hoc evaluation rather than integrated training signals.

This work addresses this gap with a simple intervention: augmenting the standard helpfulness reward with an IFEval-derived controllability signal. The combined reward takes the form:

R = $\alpha \cdot R_helpfulness + \beta \cdot R_IFEval$

where $\alpha, \beta$ control the trade-off between objectives. Discrete IFEval constraint checks are converted into continuous, differentiable rewards via sigmoid soft thresholds, enabling gradient-based optimization.

The experiments validate three predictions:

\begin{enumerate}
\item \textbf{P1 (Controllability):} Bidirectional models achieve higher held-out IFEval accuracy than helpfulness-only baselines (+3.4pp strict accuracy).
\end{enumerate}

\begin{enumerate}
\item \textbf{P2 (Helpfulness Maintenance):} At $\alpha \geq 0.8$, bidirectional models retain $\geq 95$\% of baseline AlpacaEval performance (96.4\% observed).
\end{enumerate}

\begin{enumerate}
\item \textbf{P3 (Safety Correlation):} IFEval improvements correlate with safety benchmark improvements, with +2.35pp TruthfulQA MC1 improvement over baseline maximum.
\end{enumerate}

\textbf{Contributions:}
\begin{itemize}
\item An initial empirical test of bidirectional alignment training signals, providing preliminary validation of the theoretical framework of Sun et al. (2024)
\item A methodology for repurposing rule-based evaluation metrics (IFEval) as differentiable RLHF training rewards
\item Characterization of the helpfulness-controllability trade-off at proof-of-concept scale
\end{itemize}
